% BOOK-FILE-ID: f-1d78b7d19b95
\subsubsection{丸め誤差を使う}\label{節：丸め誤差を使う}
$a_{n+1}=\lfloor\lambda a_n\rfloor$のように,無理数を掛けてから切り捨てる漸化式を考えよう.
切り捨てた分が毎回違うので,そのままでは等比型にならない.
そこで,捨てた部分に名前を付けることから始める.
\[
  \varepsilon_n=\lambda a_n-a_{n+1},
  \qquad 0\le\varepsilon_n<1
\]
と置けば,$\varepsilon_n$の範囲が分かっているので,次の整数部分を確定できることがある.
手掛かりになるのは,$\lambda^2$が$\lambda$と整数で表せることである.
ただし,誤差の範囲をきちんと示せて初めて隣接3項型に帰着できることを忘れないようにしよう.

\begin{bookpoint}
  \textbf{POINT}\par
  捨てた部分$\varepsilon_n$を置き,その範囲から次の項の整数部分を確定する.
\end{bookpoint}

\noindent \bookblocklink{example-gauss-role-5}{problem}{answer}{【例題】}\par\nobreak
\begin{enumerate}[label=(\arabic*),parsep=3pt,beginpenalty=10000,format=\bookitemlink{example-gauss-role-5}{problem}{answer}]
  \item \bookterm{sequence}{数列}を次のように定める.\bookterm{floor}{ガウス記号}を含まない\bookterm{recurrence}{漸化式}に帰着させ,\bookterm{general}{一般項}を求めよ.
    \[
      a_1=1,\qquad
      a_{n+1}=\lfloor(1+\sqrt2)a_n\rfloor
      \quad(n\ge1).
    \]
\end{enumerate}

\noindent \bookblocklink{example-gauss-role-5}{hint}{problem}{【方針】}\par\nobreak
\begin{enumerate}[label=(\arabic*),parsep=3pt,beginpenalty=10000,format=\bookitemlink{example-gauss-role-5}{hint}{problem}]
  \item $\lambda=1+\sqrt2$とおき,$\lambda^2=2\lambda+1$を使います.
    整数から正の1未満の誤差を引く形にして,\bookterm{floor}{ガウス記号}の値を確定しましょう.
\end{enumerate}

\noindent \bookblocklink{example-gauss-role-5}{answer}{problem}{【解答】}\par\nobreak
\begin{enumerate}[label=(\arabic*),parsep=3pt,beginpenalty=10000,format=\bookitemlink{example-gauss-role-5}{answer}{problem}]
  \item \[
      a_1=1,\quad a_2=2,\qquad
      a_{n+2}=2a_{n+1}+a_n-1.
    \]
    $\alpha=1+\sqrt2,\ \beta=1-\sqrt2$とおくと,
    \[
      a_n=\frac12+\frac{\alpha^n+\beta^n}{4}.
    \]
\end{enumerate}

\noindent \bookblocklink{example-gauss-role-5}{solution}{problem}{【解法】}\par\nobreak
\begin{enumerate}[label=(\arabic*),parsep=3pt,beginpenalty=10000,format=\bookitemlink{example-gauss-role-5}{solution}{problem}]
  \item $\lambda=1+\sqrt2>2$とおく.正整数$a_n$から次の項も正整数となる.
    また,$\lambda a_n$は無理数なので,
    \[
      \varepsilon_n=\lambda a_n-a_{n+1}
    \]
    は$0<\varepsilon_n<1$を満たす.
    $\lambda^2=2\lambda+1$を使うと,
    \begin{align*}
      \lambda a_{n+1}
      &=\lambda(\lambda a_n-\varepsilon_n)\\
      &=2a_{n+1}+a_n-(\lambda-2)\varepsilon_n.
    \end{align*}
    $M_n=2a_{n+1}+a_n$は整数で,$0<(\lambda-2)\varepsilon_n<1$だから,
    \[
      M_n-1<\lambda a_{n+1}<M_n.
    \]
    より$a_{n+2}=M_n-1$.\bookterm{initial}{初期条件}は$a_2=\lfloor1+\sqrt2\rfloor=2$である.

    定数項を消すため$b_n=a_n-\frac12$とおくと,
    \[
      b_{n+2}=2b_{n+1}+b_n,
      \qquad b_1=\frac12,\quad b_2=\frac32.
    \]
    \sectionref{節：隣接3項型}で学んだ\bookterm{characteristic}{特性方程式}は$x^2-2x-1=0$で,
    その解が$\alpha,\beta$である.
    $\alpha+\beta=2,\ \alpha^2+\beta^2=6$より,
    \[
      b_n=\frac{\alpha^n+\beta^n}{4}
    \]
    は初期値と漸化式を満たし,一意性から求める解となる.
    $a_n=b_n+\frac12$に戻せば【解答】を得る.
\end{enumerate}
